\documentclass[11pt]{article}
\usepackage[dvips]{graphicx}

\title{ModulusImageFilter} 
\author{Ga\"etan Lehmann}
\date{\today}

\begin{document}

\maketitle

\abstract{ModulusImageFilter is a simple image filter to compute the modulus (x \% dividend) pixel-wise.}

\section{Implementation}
ModulusImageFilter is based on UnaryFunctorImageFilter and use the c++ modulus operator (\%).

\section{Usage}
\expandafter\ifx\csname indentation\endcsname\relax%
\newlength{\indentation}\fi
\setlength{\indentation}{0.5em}
As usual, user has to include the header:
\begin{quote}
{\tt \#include} "{\tt{}itkModulusImageFilter.h}"\mbox{}\\
\end{quote}
and declare the type, instanciate the filter, and set the input image
\begin{quote}
{\bf typedef} itk::ModulusImageFilter$<$ InputType, OuputType $>$ ModulusType;\mbox{}\\
ModulusType::Pointer filter = ModulusType::New();\mbox{}\\
filter$\rightarrow$SetInput( filter$\rightarrow$GetOutput() );\mbox{}\\
\end{quote}

The dividend can be set with the SetDividend() method, and read with GetDividend(). The default value is 5.

\begin{quote}
filter$\rightarrow$SetDividend( 12 );\mbox{}\\
\end{quote}

\section{Example}
\begin{figure}
\centering
\includegraphics[width=0.5\textwidth]{Spots.eps}
\caption{The input image.}
\end{figure}

\begin{figure}
\centering
\includegraphics[width=0.5\textwidth]{distance.eps}
\caption{The (rescaled) distance map image.}
\end{figure}

\begin{figure}
\centering
\includegraphics[width=0.5\textwidth]{modulus.eps}
\caption{The (rescaled) output image.}
\end{figure}

\end{document}
